\subsection{Evitando colisões}

A inclusão de algoritmos de evasão e planejamento de movimento dinâmico é um elemento importante para garantir a segurança das missões espaciais, especialmente em ambientes complexos e dinâmicos, como órbitas terrestres, operações entre espaçonaves ou próximas a corpos celestes. Esses algoritmos permitem, como explicado por \cite{newman2024} que espaçonaves conversem entre si, detectem e evitem potenciais colisões com detritos espaciais, outros veículos espaciais ou estruturas orbitais, ajustando suas trajetórias de maneira autônoma, sem depender de comandos diretos de humanos.

O desenvolvimento e a integração desses algoritmos também são importantes para futuras missões que exijam autonomia, como operações orbitais ao redor da Lua ou de Marte. Isso permite que a espaçonave tome decisões de forma independente, reduzindo os riscos associados a atrasos na comunicação e situações inesperadas, como mudanças na dinâmica orbital ou a presença de objetos desconhecidos. Dessa forma, a evolução de sistemas de evasão e planejamento de movimento contribui diretamente para o aumento da segurança e eficiência das operações espaciais.

\subsection{Autonomia total vs. Automação avançada}

Os sistemas com automação avançada e sistemas totalmente autônomos possuem diferenças quando se trata na sua dependência de intervenção humana e também no grau de controle e decisão necessária para a situação a ser encontrada para então possam executar de maneira independente.

\subsubsection{Sistemas com Automação Avançada}

Como explicado por \cite{ibmAutomation}, são sistemas projetados para executar tarefas complexas com a mínima intervenção humana, no entanto, dependem de supervisão ou autorização para certas decisões críticas. Em um contexto espacial, sistemas com automação avançada podem executar etapas de navegação, controle e monitoramento de forma independente, mas requerem validação ou comando final de operadores humanos \cite{nasaNAC2018}, principalmente em situações de risco ou incerteza.

\subsubsection{Sistemas Totalmente Autônomos}

Os sistemas totalmente autônomos, conforme explica \cite{ibmAutomation}, operam sem qualquer intervenção humana, desde a análise do ambiente até a tomada de decisões e execução de ações. Eles são projetados para avaliar situações, identificar riscos e ajustar suas operações em tempo real, mesmo em cenários inesperados. No espaço, isso significa que uma espaçonave totalmente autônoma pode planejar e executar manobras, evitar colisões e realizar acoplamentos sem a necessidade de comunicação com operadores terrestres ou astronautas a bordo. Essa autonomia completa é particularmente útil em missões em que a latência da comunicação torna o controle em tempo real impraticável, como em operações ao redor de Marte ou em exploração em espaço profundo.

\subsubsection{Barreiras Tecnológicas}

Atualmente, a transição de sistemas automatizados para sistemas totalmente autônomos enfrenta desafios tecnológicos. Desenvolver algoritmos que possam lidar com a complexidade e imprevisibilidade de ambientes espaciais exige sensores avançados, como o LiDAR e câmeras de alta precisão, além de sistemas de processamento de dados em tempo real, o que tem sido revelado como custoso computacionalmente falando \cite{ainowComputeAI}. Esses sistemas devem ser robustos o suficiente para operar de forma confiável em condições extremas, como temperaturas extremas, radiação e ausência de gravidade. A maior preocupação é a confiabilidade, pois falhas em sistemas totalmente autônomos podem comprometer missões inteiras sem possibilidade de intervenção humana.

\subsubsection{Barreiras Regulatórias}

Além dos desafios tecnológicos, a transição para sistemas autônomos enfrenta barreiras regulatórias. A legislação espacial atual foi desenvolvida com base em operações controladas por humanos, e muitos tratados internacionais, como o Tratado do Espaço de 1967, não contemplam explicitamente a responsabilidade e os critérios de segurança para sistemas que operam sem supervisão humana. Questões sobre responsabilidade em caso de falhas, ética no uso de autonomia \cite{ainowComputeAI} e a interoperabilidade entre diferentes sistemas autônomos ainda precisam ser abordadas. Além disso, a falta de padrões técnicos globais dificulta a integração e o uso seguro desses sistemas em missões multinacionais.

\subsubsection{Interoperabilidade global}

O aumento da presença de players privados e de novos países no setor espacial tem transformado significativamente o panorama global da exploração e utilização do espaço. Empresas privadas, como SpaceX, Blue Origin, Northrop Grumman e Rocket Lab, têm expandido suas operações, oferecendo serviços de lançamento, desenvolvimento de satélites e até mesmo missões tripuladas e não tripuladas. Paralelamente, países como Emirados Árabes Unidos, Índia e Brasil \cite{senadoALADA}, assim como outros novos participantes, estão aumentando seus investimentos em tecnologias espaciais , fortalecendo suas agências espaciais e estabelecendo metas ambiciosas para explorar o espaço. Essa diversificação trouxe benefícios, como a aceleração da inovação, a redução dos custos de acesso ao espaço \cite{jones2018} e o aumento da colaboração internacional. Contudo, ela também introduziu novos desafios, principalmente no que se refere à padronização de tecnologias e protocolos.

A criação de padrões universais para interfaces \cite{nasaManualOfDocking2017} de docking é uma das prioridades no cenário espacial atual. Com a multiplicação de operadores, a falta de padronização nas interfaces de docking pode resultar em incompatibilidades que dificultam operações conjuntas, aumentam os custos e comprometem a segurança.

Essa introdução de padrões universais oferece benefícios, como:

\textbf{Interoperabilidade}: Permite que veículos espaciais de diferentes países ou empresas privadas possam se conectar sem dificuldades, promovendo a colaboração internacional.

\textbf{Redução de custos}: Evita a necessidade de desenvolvimento de tecnologias exclusivas ou adaptadores, tornando as missões mais econômicas.

\textbf{Segurança e confiabilidade}: Garante que os sistemas sejam compatíveis e testados em cenários diversos, minimizando riscos durante acoplamentos críticos.

\textbf{Incentivo à cooperação internacional}: Promove um ambiente de colaboração no espaço, unindo esforços para resolver desafios globais, como a exploração de Marte ou missões lunares.

Exemplos de iniciativas incluem o trabalho conjunto da NASA e da ESA, especialmente no contexto do programa Artemis e da Gateway \cite{nasaGateway}, uma estação espacial planejada para orbitar a Lua. O sistema de docking da Estação Espacial Internacional (ISS) já é um exemplo de cooperação, integrando tecnologias americanas e russas. Além disso, programas como o Commercial Crew Program da NASA exigem que empresas privadas, como SpaceX e Boeing, sigam padrões para garantir a compatibilidade de suas cápsulas com a ISS e futuras estações espaciais.
